\section{AegisKV Design}
\label{sec:design}

AegisKV manages reusable checkpoints in a finite shared DRAM pool and a
larger storage tier. Storage may be remote SSD or a shared filesystem;
its failure domain and performance must be specified separately. GPUs
retain their normal execution KV and workspace. There is no dedicated
Peer-HBM protection region. The design below is an implementable contract;
Section~\ref{sec:evaluation} distinguishes its analytical evaluation from
implementation evidence.

\subsection{Design Insights}
\label{sec:design-insights}

\paragraph{Recovery value depends on dependencies.}
A faster restore matters only to the extent that it advances a required
workflow result. Figure~\ref{fig:motivation-recovery-exposure} illustrates
this distinction for independent branch failures. A common worker failure
can interrupt several branches simultaneously. Their penalties cannot be
added: if two critical branches each incur 10\,s, reducing only one to
2\,s leaves their join 10\,s late. Placement must therefore evaluate the
branches affected by the same failure together.

\paragraph{Demotion and deletion have different costs.}
Moving a valid checkpoint from DRAM to storage usually adds I/O; deleting
its last surviving copy can require reconstruction. These operations
need separate values and safety checks. A large-slack agent can use
storage for recovery and still reuse its history in the next interaction.
Slack is not a reason to prohibit ordinary caching.

\paragraph{A desired placement is not recoverable state.}
Copies take time and compete with serving. A checkpoint becomes usable
only after all required objects and metadata commit. Likewise, memory is
reclaimed only after a replacement location is committed and old readers
drain. The controller must optimize the state it can actually achieve,
including the period when both copies consume capacity.

\subsection{System Overview}
\label{sec:system-overview}

The orchestrator supplies workflow IDs, dependencies, accepted inputs, and
execution progress (Figure~\ref{fig:system-overview}). The Tracker estimates
remaining durations and slack. The Router combines locality with load for
normal dispatch and selects a compatible surviving replica at failure
time. The KV Master values candidate changes to DRAM and storage residency.
Adapters expose sealed KV, resource pressure, and recovery progress;
Connectors execute transfers. A surviving directory records immutable
object identities, replica generations, and checkpoint manifests.

\begin{figure}[t]
 \centering
 \includegraphics[width=\columnwidth]{figures/storage-overview.pdf}
 \caption{AegisKV with shared DRAM and a capacity tier. Both DPs retain
 ordinary execution HBM. The Router may choose either compatible surviving
 DP; recovery reads resolve through the shared directory. Demotion commits
 the storage copy before releasing DRAM. Storage is outside the protected
 compute failure domain. Arrows show logical paths, not measured rates.}
 \label{fig:system-overview}
 \Description{A workflow supplies a Tracker and KV Master. A Router and
 surviving directory connect two inference replicas to shared DRAM and
 SSD or shared storage. There is no Peer-HBM backup region.}
\end{figure}

Objects are sealed token blocks containing one layer's complete local KV
shard. Identity includes model and weight revision, tokenizer/prefix hash,
position convention, layer, TP shard, numerical format, and token interval.
A \emph{checkpoint group} is a legal collection of these objects and
model-specific auxiliary state. Placement can move individual objects,
but valuation accounts for their recovery dependencies. Shared physical
objects consume bytes once, while their value includes all affected owners.
Dense and sparse models use the same interface; a sparse adapter must also
capture selector state. Placement never changes attention selections.

\subsection{Workflow Tracker: Remaining Time and Slack}
\label{sec:workflow-protection}

The residual DAG contains unfinished execution segments. Fork/join
boundaries split a parent so child execution is not counted again as
parent waiting. A GBDT predicts each segment's remaining duration from
role/model, elapsed time, known prompt length, and completed calls and
tools. Entry predictions mask fields unavailable before launch.
The target is remaining time directly, not total time minus elapsed time.
Train/test splits separate complete workflows and time periods.

For estimated durations $d_a$, critical-path analysis~\cite{kelley1959cpm}
computes earliest finish $E_a=d_a+\max_{b\in P(a)}E_b$, horizon
$T=\max_a E_a$, latest finish $F_a=\min_{b\in S(a)}(F_b-d_b)$, and
slack $\sigma_a=\max(0,F_a-E_a)$. Empty predecessor maxima are zero;
sinks have $F_a=T$. A request inherits its segment's margin. Multiple
requests do not each receive an independent copy of that margin.

The Tracker refreshes on arrivals, completions, DAG changes, and policy
windows, in $O(|V|+|E|)$ time after prediction. An undeclared future branch,
missing request mapping, or invalid estimate disables dependency-based
priority for that workflow. The fallback uses request recovery costs and
normal reuse, with zero slack. New DAG information invalidates the previous
snapshot. Published trajectory prediction~\cite{zhang2026heddle} motivates
using execution context; it does not establish wall-clock accuracy here.

\subsection{Recovery-aware Router}
\label{sec:recovery-router}

Normal dispatch checks compatible prefixes and foreground load, retaining
Agent affinity when useful. A protection decision does not invalidate an
existing ordinary cache entry. On a confirmed fail-stop event, the Router
fences the old worker, obtains its surviving request log and manifests,
and chooses a healthy, layout-compatible target with feasible execution
space. Shared backing removes the requirement to bind retries to a
particular Peer location. Targets can still become overloaded; both the
profile and evaluation must include readmission and background service.

Before acknowledging an output token, a surviving ordered log stores its
request ID, position, token ID, and source epoch. Prompt and accepted tool
inputs are also recoverable. Retries reconstruct the acknowledged suffix,
suppress duplicate delivery, and reject stale epochs. Exactly-once tool
side effects remain the orchestrator's responsibility. Model reload is
excluded when a compatible replica already survives; it is included when
that assumption does not hold. A GPU-process kill tests a different
failure domain from loss of a host and its local storage.

\subsection{KV Master: Online Retention and Demotion}
\label{sec:kv-master}

\subsubsection{Valuing a placement change}
\label{sec:placement-value}
Let $P$ describe committed locations and outstanding reservations, and
$e$ a possible episode within a fixed horizon $H$. Episodes include
normal next-use events and failures of named source domains. A domain
failure identifies \emph{all} its interrupted requests. The profile
returns additional segment delays $\delta(P,e)$, including checkpoint lag,
restore/reconstruction, target queues, and contended transfers. A max-plus
DAG pass evaluates workflow completion $T_w(P,e)$ using these delays.
For one isolated interrupted segment this reduces to
$\max(0,R_r(P,e)-\sigma_r)$; it is not generally the sum of such terms.

For a candidate $P'$ the controller minimizes the change
\begin{equation}
 \Delta J(P'\mid P)=A(P',P)+
 \sum_e p_e\sum_w\bigl[T_w(P',e)-T_w(P,e)\bigr].
 \label{eq:value}
\end{equation}
$A$ is the foreground delay caused by the \emph{transition}, including
writes and temporary occupancy; episode costs describe service after
that transition and exclude those same costs. All terms are seconds;
$p_e$ are probabilities or explicitly declared expected episode counts,
never raw failure rates. For a first-failure model, domain rates
$\lambda_f$ give $p_f=(\lambda_f/\Lambda)(1-e^{-\Lambda H})$ with
$\Lambda=\sum_f\lambda_f$. Multiple-failure episodes require an extended
model. Normal-use probabilities come from past Agent activations;
no future prompt or completion is visible to the controller.

Compare DRAM demotion against the actual surviving storage copy, including
writeback if absent. Compare storage deletion against the remaining valid
checkpoint or reconstruction. A group's marginal value is recomputed
against the current preview after each choice, including every owner of
a shared object. It is not copied into every block's score. Equal local
slack values can conceal a common-failure join, and the objective remains
non-additive. The greedy procedure below has no optimality or competitive
ratio guarantee.

\subsubsection{Selecting a bounded online plan}
\label{sec:placement-plan}
Newly sealed state is offered to the shared cache regardless of slack,
subject to byte and queue limits. This is an attempt, not an unlimited
storage promise. Existing compatible copies are reused first. A configured
block-aligned interval $Z$ triggers review as KV grows; arrivals, tool
handoffs, significant DAG/load changes, and capacity pressure also trigger
reviews. A review may keep every location unchanged.

A snapshot supplies legal retain, promote, demote, and discard operations.
To relieve DRAM pressure, repeatedly choose the feasible demotion with
smallest predicted increase in $J$ per byte actually releasable, breaking
ties by stable object ID. Revalue the remaining candidates after each
choice. Compare a proposed admission or promotion against this complete
victim plan, including its write cost. Accept an optional change only if
its upper cost bound is below the keep plan's lower bound by a configured
hysteresis margin. Required capacity relief can proceed without a positive
performance benefit, but still obeys safety constraints.

Storage pressure uses a separate victim set and value: only unowned
checkpoints, superseded versions, or explicitly withdrawn protection
claims can lose their last replica. With no safe victim, reject/coalesce
new writes, preserve the last valid frontier, and expose the longer replay
tail. A tool handoff releases an execution reference but does not erase
the cache or automatically withdraw a checkpoint claim. Active requests
can be demoted to storage; physical I/O leases, not permanent DRAM pins,
protect ongoing readers.

\begin{figure}[t]
 \raggedright
 \refstepcounter{designalgorithm}\label{alg:placement}
 \textbf{Algorithm \thedesignalgorithm: Bounded online tier management}
 \par\smallskip
 \begin{tabularx}{\columnwidth}{@{}r@{\quad}X@{}}
 \toprule
 1 & Snapshot DAG, epochs, locations, leases, and byte ledgers \\
 2 & Reuse committed objects; offer newly sealed dependencies \\
 3 & Form legal candidates within the configured work budget \\
 4 & While a preview exceeds a tier's capacity: \\
 5 & \quad Select least-cost legal relief per releasable byte \\
 6 & \quad Update the preview; revalue affected candidates \\
 7 & Compare keep with the complete admission/migration plan \\
 8 & Reject unsupported or unprofitable optional changes \\
 9 & Recheck versions and reserve destination/transfer space \\
 10 & Enqueue copies; retain old authoritative locations \\
 11 & On completion: validate, commit, redirect, drain, release \\
 12 & Publish actual coverage; retry after relevant state changes \\
 \bottomrule
 \end{tabularx}
 \caption*{A preview does not free memory or advance a checkpoint.}
 \Description{An online controller compares complete legal migration plans,
 reserves resources, and releases old copies only after safe commitment.}
\end{figure}

Profiles are keyed by model/layout, bytes, tail length, topology, target
load, and I/O scheduler. Each records support ranges, uncertainty, expiry,
and whether rates are measured or assumed. Missing or stale profiles
freeze discretionary promotions and use bounded LRU demotion with the
same safety rules. The prototype must measure candidate-work budgets,
profile-query overhead, and the resulting regret; unbounded exhaustive
search is only a small-instance evaluation oracle.

\subsubsection{Committing a checkpoint}
\label{sec:checkpoint}
Copies progress through pending, committed, and retiring generations.
Publication requires all required layers, TP shards, token intervals,
selector state, and a matching surviving token-log prefix. The authoritative
frontier is the latest dependency-complete manifest, not the largest
position seen in any individual copy. An incomplete new manifest leaves
the old one usable. Immutable objects can be shared across manifests.

DRAM demotion follows copy--verify--commit--redirect--drain--free. Storage
commit means the backend's documented completion/persistence guarantee;
a buffered write return is insufficient when the claimed failure model
requires durable media. Readers lease a specific location generation.
A failed copy keeps the old location, releases its failed reservation,
and triggers replanning. Duplicate completions and obsolete epochs cannot
publish a new frontier. Last-copy deletion cannot silently preserve the
old claim of recoverability.

\subsection{Adapter and Connector: Safe Concurrent Execution}
\label{sec:kv-runtime}

Each tier accounts for committed bytes, reservations, retiring buffers,
and metadata. Recovery also reserves bounded GPU workspace and staging
buffers before admission. A lower DRAM target can temporarily be below
physical occupancy while safe demotion drains; it cannot be used as an
instantaneous free-space report. Exhausted storage or transfer queues
cause bounded rejection/backpressure, not hidden allocation.

Connectors maintain separate foreground, demand-restore, and writeback
queues. Positive minimum service shares preserve progress for restore
and required demotion; unused service is borrowable. Background backup
has a configured bandwidth cap. A bandwidth share is not a bound on
foreground latency. Byte-based deficit round robin and bounded transfer
quanta prevent one long context from monopolizing a path. Normal reads
and all policies use the same scheduler in controlled comparisons.

The engine consumes chunks at supported kernel boundaries. Layer-wise
prefetch can overlap remaining replay computation, but historical
recomputation already avoided by reuse provides no overlap window.
Missing inputs stall the dependent work, not every retry. TP ranks retain
matched batch and collective order. Dense or native sparse attention
executes unchanged; batched replay requires demonstrated equivalence to
the original state evolution, otherwise the runtime replays token by token.

These rules give three invariants: pending state is never read as committed;
every published dependency resolves to a surviving compatible generation;
and a physical buffer is never reused with outstanding leases. They
establish a protocol argument, not a claim that concurrency bugs, model
parity, or fault-recovery performance have already been validated.
